\subsection{参数来源与识别边界}

为了避免把模型理解为单纯曲线拟合，本文将参数按来源分为六类。题面参数和附件统计量构成模型硬约束；官方外部数据只用于数量级校验；校准参数必须通过基准对比、消融实验和扰动检验；中长期情景参数只用于 60--180 天条件外推，不参与短期窗口的误差拟合；中长期机制强度参数则必须显式命名，并进入敏感性分析。

\begin{table}[H]
\centering
\caption{参数来源分类与使用边界}
\begin{tabularx}{\linewidth}{lX X}
\toprule
参数类别 & 代表变量 & 使用边界 \\
\midrule
题面给定参数 & 短期弹性 \(\varepsilon\)、供应中断范围、SPR 释放上限、恐慌持续时间 & 作为模型约束和传统供需基准依据，不在题面范围外任意调整 \\
附件统计量 & 基准价格 \(P_0\)、冲突窗口真实收盘价、收益率和波动率 & 作为拟合目标、误差评价和历史窗口比较依据 \\
官方外部校验量 & EIA 周度 SPR、商业库存、原油产量，JODI 多国月度石油数据，OPEC 全球供需平衡表，Caldara-Iacoviello GPR 指数，Cboe/FRED OVX 隐含波动率 & 只用于校验长期约束、数量级和风险高位事实，不替代附件真实价格 \\
机制校准参数 & 价格传导系数 \(\phi\)、风险权重 \(\beta\)、缓冲确认强度、预期修复强度、调整速度 \(v\) & 在固定机制结构内校准，并接受消融、扰动和结构稳定性检验 \\
中长期情景参数 & OPEC+ 响应、非 OPEC 供给恢复、长期弹性、制度风险强度 & 只用于条件中心路径和尾部风险外推，不称为真实未来拟合 \\
长期机制强度参数 & 冲击不确定性占比、制度风险占比、信心恢复后风险底座、过剩供给回归强度、封锁风险衰减速度 & 用于刻画长期机制形状，不作为官方物理常数，并在敏感性分析中单独扰动检验 \\
\bottomrule
\end{tabularx}
\end{table}

\FloatBarrier

其中最容易被误读的是价格传导系数 \(\phi\)。本文没有用 \(\phi\) 替换题面短期弹性 \(\varepsilon=-0.05\)：低弹性仍用于线性需求反事实基准，解释为什么静态模型会给出 278--337 美元/桶的高价压力；常弹性需求口径只作为机械上界参照。\(\phi\) 只表示在 SPR、库存、绕道运输、需求收缩和预期修复共同存在时，剩余物理缺口进入布伦特期货目标价格的传导比例。二者分别对应“无缓冲静态供需压力”和“有缓冲期货定价传导”，经济含义不同。

更具体地说，若令 \(\phi=1\)，则模型等价于假设剩余物理缺口会按题面低弹性完全、即时地映射为期货目标价格；但现实原油期货市场同时定价库存、预期、投机头寸、政策缓冲可信度和未来需求调整。Kilian（2009）、Hamilton（2009）以及 Kilian 与 Murphy（2014）均指出，油价冲击不能只由当期物理供给缺口解释，预期需求、库存和投机交易会改变现货与期货价格的传导路径。因此本文把 \(\phi\) 明确归类为行为校准参数，并在结构稳定性检验中单独报告：当 \(\phi=1\) 时，短期 RMSE 升至 26.55 美元/桶，模拟峰值升至 166.84 美元/桶，说明完全机械传导不符合附件观测；而在 \(0.5\phi\)--\(2.0\phi\) 扰动内，RMSE 仍处于 3.50--3.64 美元/桶区间，说明价格平台并非唯一参数点偶然形成。

中长期外推中的若干机制系数也需要单独说明。例如冲击不确定性占比、制度风险占比、过剩供给回归强度和封锁风险衰减速度并非 EIA、JODI 或 IEA 直接发布的物理量，而是把短期校准后的风险定价机制延伸到 60--180 天时使用的形状参数。本文不把它们写成确定事实，而是把它们显式列入敏感性分析：若这些系数在合理范围内扰动后结论仍保持中性回落、悲观尾部上行的结构，则说明中长期判断主要来自机制组合，而不是某一个隐藏常数。
